\section{面向课程实践的落地框架}

为了把讲座理念转成可执行路径，这里给出一个可用于课程项目或团队内部试验的最小实践框架。该框架不追求一步到位的 AGI，而是优先构建可闭环迭代能力。

\begin{importantbox}{四层工程栈}
\begin{enumerate}
    \item \textbf{接口层}：定义可组合 Tool API 与权限边界。
    \item \textbf{策略层}：主策略 + exploiter 策略 + evaluator 策略协同。
    \item \textbf{训练层}：IL 冷启动 + RL 强化 + match-making 调度。
    \item \textbf{评估层}：协作成功率、对抗稳健性、恢复速度、成本曲线。
\end{enumerate}
\end{importantbox}

\begin{knowledgebox}{一套可复用的实验循环}
\begin{enumerate}
    \item 采集真实任务轨迹，构建 IL 数据与失败标签。
    \item 启动主策略，使用固定安全规则做第一轮上线。
    \item 用 exploiter 代理自动挖掘脆弱点，生成 counter episodes。
    \item 以 50/50 类似准则做任务采样，避免训练信号塌缩。
    \item 周期性回归评估，分离``平均性能提高''与``尾部风险下降''。
\end{enumerate}
\end{knowledgebox}

\begin{warningbox}{课程项目里最容易被忽略的环节}
很多项目只记录最终成功率，不记录失败路径与修复时间。没有 failure replay 和修复时延统计，就很难判断系统是否真的更鲁棒，还是只是在测试集上更幸运。
\end{warningbox}

\begin{table}[H]
\centering
\caption{建议的最小指标面板}
\begin{tabular}{p{3.0cm}p{3.2cm}p{6.4cm}}
\toprule
指标组 & 代表指标 & 解释 \\
\midrule
协作效率 & 任务成功率、平均步数、工具调用次数 & 反映正常用户体验与成本 \\
对抗稳健性 & 攻击成功率、越权率、注入生效率 & 反映安全边界真实强度 \\
恢复能力 & 失败后重试成功率、平均恢复轮数 & 反映系统弹性而非一次命中 \\
分布泛化 & 新任务族成功率、跨工具迁移性能 & 反映是否过拟合既有任务 \\
\bottomrule
\end{tabular}
\end{table}

如果把上述框架与本讲观点对应，可以得到一个务实判断：\textbf{Multi-Agent 的真正价值不在于多模型并行本身，而在于它能够持续制造``有学习价值的压力''。}

\subsection{本章小结}
从实践角度看，最重要的是先搭建可闭环系统：能发现弱点、能生成反制、能验证改进。只有这样，讲座里的理论迁移才会转化为可复现的工程收益。

